\section*{Chapitre \ref{ch:m2:money-markets} --- Marchés monétaires}

\begin{solution}{exo:m2:money-markets:1}
$100\,(1 - 0.0375 \times 182/360) = 98.104167$ pour 100.
\end{solution}

\begin{solution}{exo:m2:money-markets:2}
\omterm{def:m2:money-markets:yields}{Taux d'escompte} $(100 - 99.04)/100 \times 360/91 = 3.7978\%$~; \omterm{def:m2:money-markets:yields}{taux monétaire}
$0.96/99.04 \times 360/91 = 3.8346\%$~; \omterm{def:m2:money-markets:yields}{rendement équivalent obligataire} $0.96/99.04
\times 365/91 = 3.8879\%$.
\end{solution}

\begin{solution}{exo:m2:money-markets:3}
$10\,000\,000 \times (1 + 0.04 \times 90/360) = \text{USD}~10\,100\,000$. Le taux du \omterm{def:m2:money-markets:cd}{certificat
de dépôt} est un intérêt simple sur le montant placé, actual/360~: il s'agit déjà d'un
\omterm{def:m2:money-markets:yields}{taux monétaire}, directement comparable au $m$ d'un \omterm{def:m2:money-markets:bill}{bon du Trésor}.
\end{solution}

\begin{solution}{exo:m2:money-markets:4}
En vertu de la \cref{prop:m2:money-markets:order}, $i = \frac{y}{360}\,\frac{d}{1 -
dt/360}$, produit de deux facteurs supérieurs à un, multiplié par $d$. L'écart $m - d =
\frac{d^2 t/360}{1 - dt/360}$ croît avec $t$~: à même \omterm{def:m2:money-markets:yields}{taux d'escompte}, le
rendement du bon à 26~semaines dépasse son \omterm{def:m2:money-markets:yields}{taux d'escompte} de plus que celui du bon à 4~semaines,
parce que l'escompte est calculé sur la valeur nominale et que le prix du bon le plus long est
plus nettement inférieur à sa valeur nominale.
\end{solution}

\begin{solution}{exo:m2:money-markets:5}
$P = 100\,(1 - 0.0765 \times 364/360) = 92.265000$. Avec $y = 365$~: $a =
364/730 - 0.25 = 0.248630$, $b = 364/365 = 0.997260$, $c = (92.265 -
100)/92.265 = -0.083835$, et $i = (-b + \sqrt{b^2 - 4ac})/(2a) = 8.237\%$. Un
titre à coupon de même échéance verserait un coupon au bout de six mois,
que son détenteur pourrait réinvestir~; le bon ne verse rien jusqu'à l'échéance. La
formule crédite le bon de ce réinvestissement, de sorte que son rendement se compare
à un rendement obligataire semestriel.
\end{solution}

\begin{solution}{exo:m2:money-markets:6}
Les rachats nets représentent $1.1/20 = 5.5\%$ de l'actif net, au-delà du seuil de 5~\%~: un
\omterm{def:m2:money-markets:mmf}{fonds monétaire} institutionnel prime doit prélever les frais de liquidité obligatoires auprès des
investisseurs qui rachètent. Il règle les rachats sur ses actifs liquides quotidiens, au
moins 25~\% du fonds. Un fonds gouvernemental n'est pas soumis à des frais obligatoires~; il règle
les rachats sur ses actifs liquides, qui, pour un fonds de \omterm{def:m2:money-markets:bill}{bons du Trésor} et de repo, en constituent l'essentiel.
\end{solution}

\begin{solution}{exo:m2:money-markets:7}
4~semaines 3.8642~\%, 6~semaines 3.8699~\%, 8~semaines 3.8757~\%, 13~semaines 3.8901~\%, 17
semaines 3.9018~\%, 26~semaines 3.9282~\%, 52~semaines 3.9675~\%. Une courbe d'escompte plate
est une courbe des rendements \emph{ascendante}~: la conversion ajoute davantage aux
échéances longues (l'\cref{exo:m2:money-markets:4}), et au-delà d'un semestre la
formule semestrielle prend le relais. Dix points de base de pente proviennent de la
seule convention.
\end{solution}

\begin{solution}{exo:m2:money-markets:8}
En \omterm{def:m2:money-markets:yields}{rendement équivalent obligataire}, le bon à 13~semaines rapporte 3.8901~\% et celui à 52~semaines
3.9567~\%~: la courbe \emph{monte} de 6.7~points de base. Les \omterm{def:m2:money-markets:yields}{taux d'escompte} de
différentes échéances ne sont pas des rendements comparables. Pour lire les anticipations du
marché, il faut aussi tenir compte des primes de terme et de l'offre de chaque
échéance, mais il faut d'abord comparer des grandeurs de même nature.
\end{solution}

\begin{solution}{pb:m2:money-markets:1}
\textbf{1.} Sept jours~; le placement du 31 décembre court du jeudi au lundi,
quatre jours.
\textbf{2.} $500 \times 10^6 \times 0.0395 \times 7/360 = \text{USD}~384\,028$.
\textbf{3.} $3.95\% \times 365/360 = 4.0049\%$.
\textbf{4.} $500 \times 10^6 \times [(1 + 0.0375/360)^3(1 + 0.0375 \times 4/360) -
1] = \text{USD}~364\,665$.
\textbf{5.} Le dépôt~: le trésorier prête sans collatéral à la banque. La
facilité~: aucune, mais le trésorier ne peut pas y recourir~; seuls les fonds éligibles et
les dealers peuvent.
\textbf{6.} $500 \times 10^6 \times [(1 + 0.0388/360)^3 - 1] = \text{USD}~161\,684$.
\textbf{7.} $r_T = [(1 + 0.0395 \times 7/360)/(1 + 0.0388/360)^3 - 1] \times
360/4 = 4.0009\%$.
\textbf{8.} $4.0009 - 3.88 = 12.09$ points de base.
\textbf{9.} $500 \times 10^6 \times 0.001209 \times 4/360 = \text{USD}~6\,716$.
\textbf{10.} Ses positions (titres qu'elle détient, prêts qu'elle a accordés) doivent être
financées cette nuit-là quoi qu'il arrive~; les prêteurs qui peuvent encore prendre le
risque et le bilan le facturent. Ce qu'elle ne peut pas faire à bon marché en
une seule nuit, c'est vendre ses actifs.
\textbf{11.} À la date de déclaration, les dealers et les banques ont prélevé moins de liquidités
auprès des fonds~; ceux-ci l'ont déposée à la Fed, 245~milliards de dollars de plus en un jour, et
l'ont reprise le 3~janvier.
\textbf{12.} $500 \times 10^6 \times (0.040009 - 0.0375) \times 4/360 =
\text{USD}~13\,938$.
\textbf{13.} $[(1 + 0.0388/360)^{10}(1 + 0.040009 \times 4/360) - 1] \times
360/14 = 3.9171\%$.
\textbf{14.} Les mêmes quatre jours coûteux sont moyennés sur quatorze au lieu de
sept~: le \omterm{def:m2:money-markets:turn}{passage de fin de période} ajoute 3.4 points de base au taux à deux semaines contre 6.9 au
taux à une semaine.
\textbf{15.} Ce n'est pas une coïncidence de nature~: les banques éligibles peuvent emprunter contre
des Treasuries auprès de la Fed à 4.00~\%, de sorte qu'un taux garanti très supérieur à celui-ci serait
arbitré par ceux qui disposent de collatéral et de bilan disponibles. Un
dépôt non garanti peut le dépasser, car la facilité exige du collatéral et,
comme tout emprunt, gonfle le bilan de fin d'année de l'emprunteur. Un \omterm{def:m2:money-markets:turn}{passage
de fin de période} nettement supérieur signalerait une tension, ou la réticence à être vu recourir à la facilité.
\textbf{16.} Un dépôt d'un mois à partir du 1\up{er} décembre chevauche la fin d'année~: son
taux intègre le \omterm{def:m2:money-markets:turn}{passage de fin de période}, dilué sur 31 jours, à partir du jour où il est coté.
\textbf{17.} Mesurer les bilans en moyennes sur la période plutôt que
à son dernier jour~; des réserves si abondantes qu'aucune banque n'a besoin d'emprunter pendant le
\omterm{def:m2:money-markets:turn}{passage de fin de période}~; des facilités de banque centrale qui absorbent les liquidités sans coût pour le bilan de quiconque
(ce que fait déjà la \omterm{def:m2:money-markets:rrp}{facilité de reverse repo} pour les fonds).
\textbf{18.} Gains~: aucun risque de crédit bancaire, un actif liquide. Risques et coûts~:
les bons venant à échéance juste après la fin d'année sont recherchés et rapportent en général
\emph{moins} que le dépôt~; elle doit acheter et régler le bon, et elle
perd la prime de \omterm{def:m2:money-markets:turn}{passage de fin de période} que verse la banque.
\textbf{19.} Résultat à retenir~: \emph{le \omterm{def:m2:money-markets:turn}{passage de fin de période} implicite} s'établit à 4.0009~\%, soit
une prime de 12.1 points de base par rapport aux nuits ordinaires, pour les quatre jours à partir du 31
décembre.
\textbf{20.} Parce que c'est la nuit où la taille des bilans est
mesurée, et donc la nuit où la place au bilan est rare.
\end{solution}

\begin{solution}{iq:m2:money-markets:1}
Le \omterm{def:m2:money-markets:yields}{taux d'escompte} exprime la décote par rapport à la valeur nominale~; l'investisseur paie
moins que la valeur nominale, si bien que le rendement de la somme réellement versée est plus élevé, et
le \omterm{def:m2:money-markets:yields}{rendement équivalent obligataire} du Trésor utilise aussi une année de 365 jours au lieu de 360.
Comparer avec une obligation du Trésor à deux ans sur le \omterm{def:m2:money-markets:yields}{rendement équivalent obligataire} (pour les bons dont l'échéance dépasse
un semestre, avec la formule semestrielle), puisque les rendements des obligations capitalisent
semestriellement sur une base de 365 jours.
\iqlookfor{valeur nominale contre prix payé, et 360 contre 365.}
\end{solution}

\begin{solution}{iq:m2:money-markets:2}
Le taux de la nuit où sont mesurés les bilans (fins de trimestre,
surtout la fin d'année) bondit, car les banques et les dealers réduisent leurs opérations de repo et
leurs dépôts pour cette nuit~; les prêteurs qui offrent encore de la place au bilan la font payer,
et les liquidités qui ne trouvent pas d'emprunteur vont vers les facilités de la banque centrale. La prime
apparaît dans chaque taux à terme qui couvre la date, diluée par la durée du terme.
\iqlookfor{la mesure des bilans comme cause, et la dilution}
\iqlookfor{du taux à terme.}
\end{solution}

\begin{solution}{iq:m2:money-markets:3}
Les fonds prime promettaient une valeur liquidative stable d'un dollar par action et détenaient du papier à court terme. Quand l'action d'un fonds
passait sous un dollar à cause du défaut d'un seul émetteur, les investisseurs de
chaque fonds prime avaient intérêt à demander le rachat en premier, avant que les pertes ne soient partagées.
Les fonds ont fait face aux rachats en ne renouvelant pas le papier et en le vendant~; les émetteurs qui
reposaient sur le renouvellement de leur papier ne le pouvaient plus, et le marché du papier neuf s'est arrêté. Il
a fallu une garantie du Trésor sur les actions des fonds et des facilités de la banque centrale pour
le relancer.
\iqlookfor{l'incitation à être le premier à sortir pour une créance à valeur stable sur des actifs risqués,}
\iqlookfor{et la dépendance des émetteurs au renouvellement de leurs financements.}
\end{solution}

\begin{solution}{iq:m2:money-markets:4}
La Fed est une contrepartie sans risque, de capacité illimitée à taux fixe
(jusqu'à la limite de contrepartie)~; les dealers ne peuvent pas prendre toute la liquidité, ou seulement à
un taux inférieur aux dates de déclaration~; un fonds peut avoir des plafonds de contrepartie avec
chaque dealer, ou vouloir garder ses liquidités disponibles en fin de journée, quand le
marché privé est fermé. Les 10 points de base sont le prix de ces
écarts.
\iqlookfor{capacité, plafonds de contrepartie et calendrier, et non seulement le taux.}
\end{solution}

\begin{solution}{iq:m2:money-markets:5}
$r_T = [(1 + R N/B)/(1 + r_0/B)^{N_0} - 1] B/N_T$ avec $N$ le nombre de jours du terme,
$N_0$ le nombre de jours ordinaires, $N_T$ le nombre de jours couverts par le fixing de passage (quatre si la
fin d'année tombe un jeudi précédant un jour férié). Hypothèses~: les nuits ordinaires rapportent
chacune $r_0$, capitalisées quotidiennement~; aucun autre passage sur la période~; les cotations du terme et
au jour le jour portent le même risque de crédit~; la base est de 360 jours.
\iqlookfor{les pondérations du fixing de passage, et l'énoncé des hypothèses.}
\end{solution}

\begin{solution}{iq:m2:money-markets:6}
Partir de la trajectoire du taux directeur~: les swaps de taux au jour le jour (OIS) et les
contrats à terme sur taux au jour le jour fournissent les taux au jour le jour capitalisés à chaque date.
Les bons ajoutent la composante liée à l'offre de dette publique et à la sécurité~; le repo ajoute la composante
de collatéral~; construire une courbe par classe d'instrument et modéliser leurs spreads. Intégrer
explicitement les \omterm{def:m2:money-markets:turn}{passages de fin de période}~: un saut du taux au jour le jour à terme pour chaque nuit de fin
de trimestre et de fin d'année, dimensionné à partir des cotations qui les encadrent, plutôt que de laisser
l'interpolation les étaler sur les dates voisines. Contrôler en revalorisant
chaque donnée d'entrée et la forme du taux à terme autour de chaque passage.
\iqlookfor{un modèle explicite des passages, et des courbes séparées pour les instruments}
\iqlookfor{de risques de crédit et de collatéral différents.}
\end{solution}
